\section{Local polynomials and activities}\label{sec:uv-activities}

Section~\ref{sec:uv-lattice} specified which classical observable a lattice polynomial must reproduce at small spacing.
We now construct those polynomials and introduce the form in which local interactions will be carried through the integration over scales.
The common principle is to retain the region and the coefficients of every term: this allows multiplication and summation without losing the estimates on individual observables.
An activity is a scalar function attached to a finite set of lattice sites, used to represent one term in the weight after some link variables have been integrated out.

Subsection~\ref{subsec:uv-sampling} recovers finitely many Taylor coefficients from finitely many transported plaquettes.
Subsection~\ref{subsec:uv-section-build} turns this recovery into the local representative maps of Definition~\ref{def:uv-sections}, with derivative bounds independent of the torus size.
Subsection~\ref{subsec:uv-markers} records insertions as coefficients of auxiliary commuting variables whose squares vanish.
Subsection~\ref{subsec:uv-activity-format} gives bounds for these activities and proves absolute convergence of their sums.

\subsection{Recovering Taylor coefficients}\label{subsec:uv-sampling}

The curvature of a connection~$A$ is the two-form with components
\begin{equation*}
 F_{\mu\nu}(A)=\partial_\mu A_\nu-\partial_\nu A_\mu+[A_\mu,A_\nu].
\end{equation*}
At first order in the size of~$A$, a plaquette transport records the average of the linear part~$\partial_\mu A_\nu-\partial_\nu A_\mu$ over its square. The commutator term is quadratic, so it does not enter this first variation.
Different squares give different linear measurements of the Taylor coefficients. For any fixed degree, finitely many such measurements suffice; radial gauge makes recovery of the connection from this linear curvature possible.

Fix~$d\in\N$.
Use the spaces~$\mathcal A_j,V_d$ of Definition~\ref{def:uv-labels} and the finite tuples~$J\in\mathcal J_d$ of Definition~\ref{def:uv-test-links}.
Such a tuple specifies the polynomial connection~$A^J$ and the test links~$U_{J,\varepsilon}$ defined there.
On~$\g$ use the norm~$|X|_\g=\frac{\|X\|_{\mathrm F}}{\sqrt m}$, and let~$P_\g$ be real Frobenius orthogonal projection onto~$\g$.
The following matrix-valued function maps a plaquette matrix to the Lie algebra; its derivative at the identity is the identity map, so it retains the linear curvature term just described. Here~$u\in G$, and~$\chi$ is the distance from the identity in~\eqref{eq:uv-chord}:
\begin{equation}\label{eq:uv-curvature}
 \mathcal F(u)=P_\g\!\left(\frac{u-u^{-1}}{2}\right),\qquad
 |\mathcal F(u)|_\g\le\chi(u),\qquad
 \left.\frac{d}{ds}\mathcal F(e^{sX})\right|_{s=0}=X\quad(X\in\g)\,.
\end{equation}
Projection commutes with conjugation by~$G$, since conjugation is orthogonal and preserves~$\g$.
The bound follows by writing~$u-u^{-1}=(u-1)-(u^{-1}-1)$ and using~\eqref{eq:uv-chord-square}.
The derivative follows from~$e^{sX}-e^{-sX}=2sX+O(s^3)$.
Also~$\mathcal F(u^{-1})=-\mathcal F(u)$.

For a plaquette~$p$ based at~$x\in\Z^4$ and a walk~$\omega$ from~$0$ to~$x$, transport its curvature back to the origin:
\begin{equation}\label{eq:uv-transported-curvature}
 Q_{\omega,p}(U)=\Ad\bigl(U(\omega)\bigr)\mathcal F(U_p),\qquad
 Q_{\omega,p}(U^\gamma)=\Ad\bigl(\gamma(0)\bigr)Q_{\omega,p}(U)\,.
\end{equation}
Here~$U(\omega)$ is the ordered walk transport of Definition~\ref{def:uv-walks}, $U_p$ is the plaquette matrix of Definition~\ref{def:plaquette}, and~$\Ad(g)X=gXg^{-1}$ as in Definition~\ref{def:uv-labels}.
The second identity uses the gauge transformation of Definition~\ref{def:gauge-transform} and~\eqref{eq:uv-gauge}.
All these quantities now transform at the same site, so invariant polynomials in them will be gauge invariant.
They read only links in the walk and plaquette, with both endpoints included.

The next lemma supplies a local analytic recovery map.
It uses several plaquettes because their combined measurements have enough independent components to recover all Taylor coefficients through degree~$d$.

\begin{lemma}[Finite sampling]\label{lem:uv-sampling}
There is a finite family~$(\omega_a,p_a)_{a=1}^{M_d}$ for which
\begin{equation}\label{eq:uv-sampler}
 \mathcal Q_d(J)=\bigl(Q_{\omega_a,p_a}(U_{J,1})\bigr)_{a=1}^{M_d},
 \qquad \mathcal Q_d:\mathcal J_d\longrightarrow\g^{M_d},
\end{equation}
has injective derivative at~$0$.
There is a real analytic map~$R_d:\g^{M_d}\to\mathcal J_d$, defined near~$0$, satisfying
\begin{equation}\label{eq:uv-recovery}
 R_d(0)=0,\qquad R_d\mathcal Q_d=\mathrm{id},\qquad
 R_d(\Ad(g)q)=\Ad(g)R_d(q)\quad(g\in G)\,.
\end{equation}
The last two equalities hold near the origin, for each fixed~$g\in G$; the test links are from Definition~\ref{def:uv-test-links}, and~$\Ad(g)$ acts on every component of a tuple.
\end{lemma}

\begin{proof}
The recovery has three steps: all square measurements together detect every jet, finite dimensionality reduces them to finitely many measurements, and the inverse function theorem recovers the jet from their nonlinear values.
The two properties to establish are
\begin{equation*}
 \mathcal Q_d(J)=\bigl(Q_{\omega_a,p_a}(U_{J,1})\bigr)_{a=1}^{M_d},
 \qquad R_d\mathcal Q_d(J)=J\ \text{ near }0.
\end{equation*}

\begin{enumerate}
\item \emph{All square measurements detect the jet.}

Recall that~$A_\mu^J=\sum_{j+2\le d}A_{j,\mu}$ is the truncated connection of Definition~\ref{def:uv-test-links}.
Vary its size by~$\lambda\in\R$ in the transport equation~\eqref{eq:uv-link-ode}, keeping the edge fixed.
At~$\lambda=0$ the solution is~$1$, so its first variation is the integral along that edge:
\begin{equation*}
 \left.\frac{d}{d\lambda}U_{\lambda J,1}(x,\mu)\right|_{\lambda=0}
 =\int_0^1 A_\mu^J(x+te_\mu)\,dt\,.
\end{equation*}
Since~$\mathcal F(1)=0$, differentiating the conjugating transport in~\eqref{eq:uv-transported-curvature} contributes zero.
Consequently, for the unit square~$p=(x;\mu,\nu)$,
\begin{equation}\label{eq:uv-linear-square}
 D\bigl(Q_{\omega,p}(U_{J,1})\bigr)_{J=0}[J]
 =\int_0^1\!\int_0^1
   (\partial_\mu A_\nu^J-\partial_\nu A_\mu^J)(x+se_\mu+te_\nu)\,ds\,dt\,.
\end{equation}
This is independent of~$\omega$.

Suppose all these integrals vanish for integer~$x$.
Each is a polynomial in~$x$; a polynomial vanishing on~$\Z^4$ vanishes everywhere, by successive one-variable polynomial identities.
For the highest homogeneous term~$A_j$, the highest part of that integral is~$\partial_\mu A_{j,\nu}-\partial_\nu A_{j,\mu}$.
It must therefore vanish.
Differentiating the radial identity~\eqref{eq:uv-radial} and using homogeneity gives
\begin{equation}\label{eq:uv-radial-recovery}
 (j+2)A_{j,\mu}(z)
 =\sum_\nu z^\nu(\partial_\nu A_{j,\mu}-\partial_\mu A_{j,\nu})(z)\,.
\end{equation}
Thus~$A_j=0$, and descending in degree proves~$J=0$.

\item \emph{Retain finitely many measurements.}
The kernels of the linear measurements have zero intersection.
While their intersection is nonzero, another measurement decreases its dimension, so finitely many suffice.
For~$d<2$ the source space is zero and the assertion is immediate.

\item \emph{Invert the finite nonlinear map.}
The map~$\mathcal Q_d$ is analytic near zero by the transport series~\eqref{eq:uv-transport-series}.
Set~$B=D\mathcal Q_d(0)$.
Conjugating the connection by a constant~$g\in G$ conjugates its test links, by uniqueness of~\eqref{eq:uv-link-ode}.
Equation~\eqref{eq:uv-transported-curvature} and differentiation therefore give
\begin{equation*}
 \mathcal Q_d(\Ad(g)J)=\Ad(g)\mathcal Q_d(J),\qquad
 B\Ad(g)=\Ad(g)B\,.
\end{equation*}
The adjoint action on~$\g^{M_d}$ is orthogonal and preserves the range~$\operatorname{ran}B=\{BJ:J\in\mathcal J_d\}$ of~$B$ in both directions.
Let~$\Pi$ be the orthogonal projection onto this range; it commutes with that action.
Here~$B^{-1}$ denotes the inverse from~$\operatorname{ran}B$ to~$\mathcal J_d$.
The linear map~$C=B^{-1}\Pi$ satisfies
\begin{equation*}
 CB=\mathrm{id},\qquad C\Ad(g)=\Ad(g)C,\qquad
 D(C\mathcal Q_d)(0)=\mathrm{id}\,.
\end{equation*}
Since~$\mathcal Q_d(0)=0$, the analytic inverse function theorem gives a local inverse~$g_d$ of~$C\mathcal Q_d$ fixing~$0$.
Then~$R_d=g_dC$ has the first two properties in~\eqref{eq:uv-recovery}.
The last identity in~\eqref{eq:uv-recovery} follows from uniqueness of the local inverse, for each fixed~$g\in G$.
\end{enumerate}
\end{proof}

\subsection{Constructing the local polynomials}\label{subsec:uv-section-build}

The recovery map is analytic only near the zero connection.
Its Taylor polynomial, however, can be evaluated on every link configuration.
Composing this polynomial with the classical observable polynomials~$v\in V_d$ of Definition~\ref{def:uv-labels} gives the required local representatives.
A finite average restores the signed coordinate symmetries exactly.

We measure derivatives by varying each link matrix within~$G$.
The permitted size of a variation grows with its distance from the insertion site, so we first specify that distance.
On the lattice~$\Lambda_\ell$ of Definition~\ref{def:lattice}, periodic maximum distance is
\begin{equation}\label{eq:uv-periodic-distance}
 d_\ell(x,y)=\min_{z\in\ell\Z^4}\|\widetilde x-\widetilde y+z\|_\infty\,,
\end{equation}
where tildes denote arbitrary integer representatives and~$\|z\|_\infty=\max_\mu|z^\mu|$.
For~$0<t\le1$ and a set of sites~$X\subseteq\Lambda_\ell$, define the bad sites and the configurations good on~$X$ by
\begin{equation}\label{eq:uv-good-configurations}
 \begin{aligned}
 \mathcal B_t(U)&=\{x:\text{some plaquette $p$ containing $x$ has $\chi(U_p)>t$}\},\\
 \mathcal G_t(X)&=\{U:X\cap\mathcal B_t(U)=\varnothing\}\,.
 \end{aligned}
\end{equation}
Every plaquette touching~$X$ is tested, and equality~$\chi(U_p)=t$ belongs to the good set.
The bad-site indicators are Borel and gauge invariant, by continuity and conjugation invariance of~$\chi$.

Fix an insertion site~$o\in\Lambda_\ell$ and integers~$k,q\in\N$.
Here~$k$ is the largest derivative order and~$q$ controls the permitted growth with distance.
A direction~$A$ assigns an element of~$\g$ to each positive link~$e$.
The path~$U(e)\mapsto e^{sA(e)}U(e)$ changes each matrix by multiplication on the left and remains in~$G$ for real~$s$ (Definition~\ref{def:gauge}).
Differentiate along this path and impose the bound
\begin{equation}\label{eq:uv-left-derivative}
 D_AF(U)=\left.\frac{d}{ds}F\bigl((e^{sA(e)}U(e))_e\bigr)\right|_{s=0},\qquad
 |A(e)|_\g\le t\bigl(1+d_\ell(e^-,o)\bigr)^q\,.
\end{equation}
Here~$e^-$ is the initial site of the link.
Write~$\operatorname{Dir}_{t,q}(o)$ for the directions satisfying the bound in~\eqref{eq:uv-left-derivative}, and define
\begin{equation}\label{eq:uv-activity-seminorm}
 \|F\|_{t,k,q;o,X}
 =\sup_{\substack{U\in\mathcal G_t(X),\ 0\le r\le k\\
                   A_1,\ldots,A_r\in\operatorname{Dir}_{t,q}(o)}}
       |D_{A_1}\cdots D_{A_r}F(U)|\,.
\end{equation}
For~$r=0$, the expression under the supremum is~$|F(U)|$.
This is a seminorm on functions with these derivatives; no value outside~$\mathcal G_t(X)$ is measured.

The small factor in the following bound counts the number of classical factors, independently of their weights.
The local representative maps are chosen before the derivative order, the threshold and the torus.

\begin{proposition}[Normalised local representative maps]\label{prop:uv-sections}
For the bases~$b_{d,i}$ and factor numbers~$f_{d,i}$ of Definition~\ref{def:uv-bases}, one can choose local representative maps~$P_d$ satisfying Definition~\ref{def:uv-sections}, simultaneously for all~$d$.
Each~$P_d[b_{d,i}]$ is a polynomial in real coordinates of finitely many quantities~\eqref{eq:uv-transported-curvature}, with every monomial containing at least~$f_{d,i}$ factors.
For every~$k,q\in\N$ with~$k\ge6$, $q\ge2$, there is~$A_{d,k,q}\ge0$ such that
\begin{equation}\label{eq:uv-section-bound}
 \bigl\|P_d[b_{d,i}](o;\cdot)\bigr\|_{t,k,q;o,\,o+[\Sigma_d]_\ell}
 \le A_{d,k,q}t^{f_{d,i}}
 \qquad(\ell\ge1,\ o\in\Lambda_\ell,\ 0<t\le1)\,.
\end{equation}
\end{proposition}

\begin{proof}
The construction has three parts: approximate the inverse to the required Taylor order, average over signed coordinate changes, and estimate each product of transported plaquettes on the good set.
The two key controls are
\begin{equation*}
 \left.\frac{d^r}{d\varepsilon^r}R_d^{[d]}\!\left(\mathcal Q_d(\operatorname{sc}_\varepsilon J)\right)\right|_{\varepsilon=0}
 =\left.\frac{d^r}{d\varepsilon^r}\operatorname{sc}_\varepsilon J\right|_{\varepsilon=0}\quad(0\le r\le d),
 \qquad \|P_d[b_{d,i}]\|_{t,k,q;o,\,o+[\Sigma_d]_\ell}
 \le A_{d,k,q}t^{f_{d,i}}.
\end{equation*}
By~\eqref{eq:uv-recovery}, the composition~$R_d^{[d]}\mathcal Q_d(J)$ has the same Taylor coefficients as~$J$ through order~$d$ at~$J=0$.

\begin{enumerate}
\item \emph{Match the Taylor coefficients and retain factors.}

Let~$R_d^{[d]}$ be the ordinary Taylor polynomial at~$0$ of the recovery map~$R_d$ in Lemma~\ref{lem:uv-sampling}, through degree~$d$.
Apply it to the finite tuple of transported plaquettes chosen in that lemma and set
\begin{equation}\label{eq:uv-unsymmetrized-section}
 P_d^0[v](U)=v\!\left(R_d^{[d]}\bigl((Q_{\omega_a,p_a}(U))_a\bigr)\right)\,.
\end{equation}
The polynomial~$R_d^{[d]}$ has zero constant term, real coefficients, and commutes with conjugation by~$G$, by differentiating~\eqref{eq:uv-recovery}.
Hence~$P_d^0$ is gauge invariant, respects coefficient conjugation, and sends~$1$ to~$1$.
Inverses of unitary matrices are conjugate transposes, so~\eqref{eq:uv-unsymmetrized-section} is a polynomial function of link entries and their conjugates.
Substituting polynomials with zero constant term into a classical observable polynomial with~$f_{d,i}$ factors leaves at least that many transported-curvature factors in every monomial.
Only the recovery map~$R_d$ is replaced by its degree-~$d$ Taylor polynomial. After composition with~$v$, we keep the full resulting polynomial in the link variables, so it is defined on every link configuration.

For the scaling of jets in Definition~\ref{def:uv-test-links}, homogeneity in the transport equation and uniqueness give
\begin{equation}\label{eq:uv-scaled-jet}
 U_{J,\varepsilon}=U_{\operatorname{sc}_\varepsilon J,1}\,.
\end{equation}
The compositions~$R_d^{[d]}\mathcal Q_d$ and~$R_d\mathcal Q_d$ have equal Taylor coefficients through degree~$d$ at zero.
Composing with~$\operatorname{sc}_\varepsilon J$ and then~$v$ proves~\eqref{eq:uv-section-taylor} for~$P_d^0$.

\item \emph{Restore the exact signed symmetries.}
Average over the group of Definition~\ref{def:uv-symmetries}:
\begin{equation}\label{eq:uv-section-average}
 P_d[v](U)=\frac{1}{|H_4|}\sum_{h\in H_4}P_d^0[h^{-1}\cdot v](h^*U)\,.
\end{equation}
Changing the summation index gives the signed coordinate identity in~\eqref{eq:uv-section-symmetry}; the unit and reality identities pass to the average as well.
For the Taylor identity, the test links satisfy
\begin{equation}\label{eq:uv-holonomy-covariance}
 h^*U_{J,\varepsilon}=U_{T_{h^{-1}}J,\varepsilon}\qquad(h\in H_4)\,.
\end{equation}
This follows by substitution in~\eqref{eq:uv-link-ode} and uniqueness; on a reversed link the solution is~$u(1)^{-1}u(1-t)$.
In each summand of~\eqref{eq:uv-section-average}, the inverse action on polynomials from Definition~\ref{def:uv-actions} cancels the change in the test jet, preserving~\eqref{eq:uv-section-taylor}.
A transformed transported plaquette is another transported plaquette, with a possible sign from reversal of its circulation.
Since the action on polynomials preserves factor number, the lower factor bound survives averaging.
Enclose all walks and plaquettes occurring in this finite sum in one centred integer cube~$\Sigma_d$.
This gives every condition in Definition~\ref{def:uv-sections}.

\item \emph{Bound derivatives on the good set.}
For the bound, each transported curvature~$Q$ has value at most~$t$ on the good set.
All links occurring in its fixed walks and plaquettes have maximum distance at most a fixed~$R$ from the origin on~$\Z^4$.
After periodic placement their distance is still at most~$R$, so the allowed direction there has size at most~$t(1+R)^q$.
Each differentiation of a link product introduces one such direction.
For~$1\le r\le k$, this gives
\begin{equation}\label{eq:uv-factor-derivatives}
 |Q(U)|_\g\le t,\qquad
 |D_{A_1}\cdots D_{A_r}Q(U)|_\g\le C_{r,q}t^r\le C_{r,q}t\,.
\end{equation}
The constant includes the fixed matrix projection and finitely many path factors.
Leibniz's rule therefore gives~$Ct^s$ for every derivative through order~$k$ of a monomial with~$s$ factors.
Sum its finitely many coefficients and use~$s\ge f_{d,i}$ and~$t\le1$ to obtain~\eqref{eq:uv-section-bound}.
Repeated links on a small torus affect only the finite differentiation multiplicities, so the constant is independent of~$\ell$.
\end{enumerate}
\end{proof}

These polynomial functions are smooth along the link variations~\eqref{eq:uv-left-derivative} at every configuration; the estimate applies on the closed good set, including its boundary.
The derivatives are taken on the polynomial itself, before multiplying by any indicator of the good set.
When regular terms are defined in Subsection~\ref{subsec:uv-activity-format}, their derivatives are likewise taken before multiplication by the good-set indicator.

\subsection{Products of insertions}\label{subsec:uv-markers}

To record one observable in a weighted product, introduce a commuting formal variable~$\zeta$ with~$\zeta^2=0$.
The coefficient of~$\zeta$ then records exactly one insertion.
Several such variables record finite lists of insertions, and restrictions on allowed lists can be imposed in the same algebra.

Let~$V$ be a finite set indexing the formal variables, and let~$\mathcal R\subseteq2^V$ be the collection of allowed sets of insertions.
Assume that~$\mathcal R$ is closed under taking subsets and contains~$\varnothing$ and every singleton.
For~$I\in\mathcal R$, write~$\zeta_I=\prod_{v\in I}\zeta_v$.
Define the algebra and its weighted coefficient norm by
\begin{equation}\label{eq:uv-marker-algebra}
 \begin{gathered}
 \mathcal A_{\mathcal R}=\bigoplus_{I\in\mathcal R}\C\zeta_I,\qquad
 a=\sum_{I\in\mathcal R}a_I\zeta_I,\quad 1=\zeta_\varnothing,\\
 \zeta_I\zeta_J=
 \begin{cases}\zeta_{I\cup J},&I\cap J=\varnothing,\ I\cup J\in\mathcal R,\\0,&\text{otherwise},\end{cases}\\
 \|a\|_\omega=\sum_{I\in\mathcal R}\omega^I|a_I|,\qquad
 \omega^I=\prod_{v\in I}\omega_v,\quad \omega_v>0\,.
 \end{gathered}
\end{equation}
The variables commute; their squares, and products excluded by~$\mathcal R$, are zero. For three factors, either both orders of multiplication give zero, or both give the monomial indexed by the union of their pairwise disjoint variable sets; hence multiplication is associative.

The weights multiply exactly when disjoint insertions are combined.
Thus the estimates below have no extra factor depending on the number of formal variables.

\begin{proposition}[Square-zero products]\label{prop:uv-marker}
For~$a,b\in\mathcal A_{\mathcal R}$ of~\eqref{eq:uv-marker-algebra},
\begin{equation}\label{eq:uv-marker-norm}
 \|1\|_\omega=1,\qquad \|ab\|_\omega\le\|a\|_\omega\|b\|_\omega,\qquad
 \left\|\sum_{r=0}^{|V|}\frac{a^r}{r!}\right\|_\omega\le e^{\|a\|_\omega}\,.
\end{equation}
If~$a_\varnothing=b_\varnothing=0$, then~$a^{|V|+1}=0$ and the finite exponential~$\operatorname{Exp}(a)=\sum_{r=0}^{|V|}\frac{a^r}{r!}$ satisfies
\begin{equation}\label{eq:uv-marker-exponential}
 \operatorname{Exp}(a+b)=\operatorname{Exp}(a)\operatorname{Exp}(b),\qquad
 \operatorname{Exp}(-a)\operatorname{Exp}(a)=1\,.
\end{equation}
\end{proposition}

\begin{proof}
The norm estimate follows because only disjoint variable sets multiply nontrivially, and their weights multiply.
The exponential identities follow from nilpotence and the finite binomial formula.
We prove these as follows.
\begin{enumerate}
\item \emph{Product bound.} A nonzero coefficient in~$ab$ comes from disjoint sets~$I,J$, with~$\omega^{I\cup J}=\omega^I\omega^J$; hence
\begin{equation*}
 \|ab\|_\omega
 \le\sum_{\substack{I,J\in\mathcal R\\I\cap J=\varnothing,\ I\cup J\in\mathcal R}}
       \omega^I\omega^J|a_I||b_J|
 \le\left(\sum_I\omega^I|a_I|\right)\left(\sum_J\omega^J|b_J|\right)\,.
\end{equation*}

\item \emph{Norm of the finite exponential.} Since~$\|1\|_\omega=1$, submultiplicativity and the triangle inequality give
\begin{equation*}
 \|\operatorname{Exp}(a)\|_\omega
 \le\sum_{r=0}^{|V|}\frac{\|a\|_\omega^r}{r!}
 \le e^{\|a\|_\omega}.
\end{equation*}

\item \emph{Exponential identities.}
If the constant coefficient vanishes, every factor of~$a$ uses a nonempty set of variables.
A product of more than~$|V|$ such factors is zero.
Thus~$a^{|V|+1}=b^{|V|+1}=0$, and the binomial formula yields
\begin{equation*}
 \operatorname{Exp}(a+b)=\operatorname{Exp}(a)\operatorname{Exp}(b),
 \qquad \operatorname{Exp}(-a)\operatorname{Exp}(a)=1.
\end{equation*}
\end{enumerate}
\end{proof}

The same norm justifies countable summation: if~$\|a_j\|_\omega\le c_j$ and~$\sum_jc_j<\infty$, the partial sums are Cauchy and~$\|\sum_ja_j\|_\omega\le\sum_jc_j$.
Completeness follows from finite dimensionality.
For one variable, let~$\rho$ be a scalar weight and~$I$ an observable, viewed as functions of the link configuration.
Multiplication by~$1+\zeta I$ records the term containing one copy of~$I$ while leaving the weight attached.
For a finite index set~$\mathcal I$, let~$K_j$ be scalar factors and~$H_j$ the factors with an insertion, for~$j\in\mathcal I$.
The insertion rule is then
\begin{equation}\label{eq:uv-one-insertion}
 [\zeta]\,\rho(1+\zeta I)=\rho I,\qquad
 [\zeta]\prod_{j\in\mathcal I}(K_j+\zeta H_j)
 =\sum_{j\in\mathcal I}H_j\prod_{i\in\mathcal I\setminus\{j\}}K_i\,.
\end{equation}
Here~$[\zeta]$ means the coefficient of~$\zeta$.
For any measure~$\mu$ for which these functions are integrable, linearity of integration gives the corresponding identity
\begin{equation*}
 \int [\zeta]\rho(1+\zeta I)\,d\mu
 =\int\rho I\,d\mu
 =[\zeta]\int\rho(1+\zeta I)\,d\mu.
\end{equation*}
For the finite product in~\eqref{eq:uv-one-insertion}, the same fact is the displayed identity
\begin{equation*}
 \int [\zeta]\prod_{j\in\mathcal I}(K_j+\zeta H_j)\,d\mu
 =\sum_{j\in\mathcal I}\int H_j\prod_{i\in\mathcal I\setminus\{j\}}K_i\,d\mu
 =[\zeta]\int\prod_{j\in\mathcal I}(K_j+\zeta H_j)\,d\mu,
\end{equation*}
provided the terms are integrable.
Thus the coefficient of an insertion in an expectation is obtained by integrating the product of that observable and its weight.

\subsection{Local interactions}\label{subsec:uv-activity-format}

When links at one lattice spacing are integrated out, the resulting function of the remaining links is the effective weight.
We seek to write it as a sum of scalar terms, each assigned to a finite set of sites; different indices may refer to the same set.
For a region~$X$, the good-set condition~$\mathcal G_t(X)$ from~\eqref{eq:uv-good-configurations} means that every plaquette touching~$X$ has~$\chi(U_p)\le t$.
Activities controlled by derivatives on this good set are the regular terms. Terms supported on configurations with large plaquette distances retain the sites where~$\chi(U_p)>t$; the formula below calls these bad sites.

Use the lattice of Definition~\ref{def:lattice}, the threshold~$0<t\le1$, the derivative order~$k\ge6$ and distance exponent~$q\ge2$ from~\eqref{eq:uv-activity-seminorm}.
The support weight counts the balls in a smallest connected cover: each ball contributes a factor~$e^\alpha$, so a larger or more separated region receives a larger weight.

\begin{definition}[Connected ball covers]\label{def:uv-covers}
Fix an integer radius~$B\ge1$ and~$\alpha>0$.
A ball of radius~$B$ about~$a$ is~$\mathbb B_{\ell,B}(a)=\{x\in\Lambda_\ell:d_\ell(x,a)\le B\}$.
A set of sites is connected if any two of its sites can be joined within the set by a finite sequence of steps of maximum distance at most one.
For~$X\subseteq\Lambda_\ell$, define
\begin{equation}\label{eq:uv-cover-number}
 \begin{aligned}
 \mathfrak c_B(X)&=\min_{\substack{D\subseteq\Lambda_\ell\\
                    X\subseteq\bigcup_{a\in D}\mathbb B_{\ell,B}(a)\\
                    \bigcup_{a\in D}\mathbb B_{\ell,B}(a)\ \mathrm{connected}}}|D|\,.
 \end{aligned}
\end{equation}
Thus~$\mathfrak c_B(\varnothing)=0$; the support weight is~$e^{\alpha\mathfrak c_B(X)}$.
\end{definition}

The minimum exists: taking every site as a centre covers the torus, which is connected by coordinate steps.

We call the activities controlled by derivatives on good configurations the regular terms.
Fix finite sets of sites~$X_i\subseteq\Lambda_\ell$, sites~$o_i\in X_i$, and real Borel gauge-invariant functions~$E_i$ supported on the links inside~$X_i$ in Definition~\ref{def:observables}.
They have the derivatives measured by~\eqref{eq:uv-left-derivative} on the good set~$\mathcal G_t(X_i)$ of~\eqref{eq:uv-good-configurations}.
Their contribution is~$\mathbf1_{\mathcal G_t(X_i)}E_i$.
For the other terms, fix finite site sets~$Y_j\subseteq\Lambda_\ell$. Choose a nonempty set~$I_j\subseteq Y_j$ of sites that the term requires to be bad, and an optional extra condition~$\mathcal E_j$ on configurations; if no extra condition is imposed, take~$\mathcal E_j=\mathcal U_\ell$. Require~$\mathcal E_j$ to be Borel and gauge invariant. For each~$j$, choose a real Borel gauge-invariant function~$K_j$ that vanishes unless all sites of~$I_j$ and no other sites of~$Y_j$ are bad and~$\mathcal E_j$ holds:
\begin{equation}\label{eq:uv-forced-eligibility}
 K_j(U)\ne0\ \Longrightarrow\ Y_j\cap\mathcal B_t(U)=I_j
 \quad\text{and}\quad U\in\mathcal E_j\,.
\end{equation}
Here the internal links of~$Y_j$ are those with both endpoints in~$Y_j$, as in the link-support convention of Definition~\ref{def:observables}. Both~$K_j$ and~$\mathcal E_j$ depend only on those links and the bad-site indicators in~$Y_j$.
The fixed sets and extra conditions are retained when an insertion changes the coefficients.

Evaluating the bad-site indicators of~\eqref{eq:uv-good-configurations}, either for a regular term or for the conditions~\eqref{eq:uv-forced-eligibility}, may require extra links.
For a set of sites~$Y$, write~$\operatorname{vertices}(p)$ for the four corner sites of a plaquette~$p$.
All such links have endpoints in
\begin{equation}\label{eq:uv-inspected-support}
 \widehat Y=Y\cup\bigcup_{\operatorname{vertices}(p)\cap Y\ne\varnothing}
                 \operatorname{vertices}(p)\,.
\end{equation}
The two activity bounds count terms differently: regular terms are summed at their chosen site, and the other terms at every site in their support.
Precisely, bounds~$\delta,w\ge0$ mean that there exist nonnegative majorants~$e_i,k_j$ (upper bounds for the individual activities) with
\begin{equation}\label{eq:uv-activity-bounds}
 \begin{aligned}
 \|E_i\|_{t,k,q;o_i,X_i}&\le e_i,& \sup_U|K_j(U)|&\le k_j,\\
 \sum_{i:o_i=x}e^{\alpha\mathfrak c_B(X_i)}e_i&\le\delta,&
 \sum_{j:x\in Y_j}e^{\alpha\mathfrak c_B(Y_j)}k_j&\le w
 \quad(x\in\Lambda_\ell)\,.
 \end{aligned}
\end{equation}
The last two sums are convergent series of nonnegative terms.
Indices range over~$\N$ and zero functions may be added.

Let~$c\in\R$ be a chosen constant per site; it contributes the overall factor~$e^{c|\Lambda_\ell|}$ to the weight. Let~$S$ be the part of the action kept separate from the activities, assumed bounded, real, Borel and gauge invariant; write~$\|S\|_\infty=\sup_U|S(U)|$.
For a finite selection~$\mathcal I\subset\N$ of indices, multiply the factors~$K_j$.
The sum~$\mathcal Z_K$ retains only disjoint regions whose specified bad sites cover the entire measured bad set; $E_t$ sums the regular contributions:
\begin{equation}\label{eq:uv-activity-weight}
 \begin{aligned}
 E_t(U)&=\sum_i\mathbf1_{\mathcal G_t(X_i)}(U)E_i(U),\\
 \mathcal Z_K(U)&=\sum_{\substack{\mathcal I\subset\N\ \mathrm{finite}\\
              Y_j\ (j\in\mathcal I)\ \mathrm{pairwise\ disjoint}\\
              \bigcup_{j\in\mathcal I}I_j=\mathcal B_t(U)}}\ \prod_{j\in\mathcal I}K_j(U),\\
 \rho(U)&=\exp\!\bigl(c|\Lambda_\ell|-S(U)-E_t(U)\bigr)\mathcal Z_K(U)\,.
 \end{aligned}
\end{equation}
The empty selection contributes~$1$ precisely when~$\mathcal B_t(U)=\varnothing$.
The sum requires the regions~$Y_j$ to be disjoint, but their inspected sets~$\widehat Y_j$ may overlap: those larger sets only record the extra links needed to decide which sites are bad.
The function~$\rho$ is an unnormalised scalar factor intended to be integrated against product Haar; positivity and probability normalisation are not asserted here.

These formulas are useful only if their infinite sums converge absolutely.
The per-site bounds give this convergence at each finite volume. To identify the represented weight with one obtained from the Wilson measure, a later step must:
\begin{enumerate}
\item integrate links at one spacing;
\item collect the resulting terms by finite site region; and
\item verify the prescribed values and Taylor terms of total degree two at the identity, together with the activity bounds.
\end{enumerate}
Here the identity configuration is~$U(x,\mu)=1$ for every positive link. Nearby links are written~$U(x,\mu)=\exp(X_{x,\mu})$ with~$X_{x,\mu}\in\g$ near zero, so the degree-two terms are the homogeneous Taylor terms of total degree two in these coordinates. This construction is not supplied here.

\begin{proposition}[Convergence of the activity sums]\label{prop:uv-activity-sums}
Under~\eqref{eq:uv-forced-eligibility}--\eqref{eq:uv-activity-bounds}, both series in~\eqref{eq:uv-activity-weight} converge absolutely and uniformly in~$U$.
They are Borel and gauge invariant, and
\begin{equation}\label{eq:uv-activity-global-bounds}
 \begin{aligned}
 |E_t(U)|&\le|\Lambda_\ell|\delta,&
 |\mathcal Z_K(U)|&\le e^{|\Lambda_\ell|w},\\
 |\rho(U)|&\le
 \exp\!\bigl(c|\Lambda_\ell|+\|S\|_\infty+|\Lambda_\ell|(\delta+w)\bigr)\,.
 \end{aligned}
\end{equation}
In particular~$\rho$ is integrable against the product Haar measure~$\eta_\ell$ of Definition~\ref{def:wilson-law}.
If~$\mathcal B_t(U)=\varnothing$, then~$\mathcal Z_K(U)=1$.
\end{proposition}

\begin{proof}
At fixed finite volume, the per-site hypotheses give summable bounds for individual terms; the sum over compatible selections is then bounded by an unrestricted product.
The argument has three steps.

\begin{enumerate}
\item \emph{Sum individual activities.}
Since each regular term has one chosen site and every~$Y_j$ is nonempty, the weights in~\eqref{eq:uv-activity-bounds} are at least one and give
\begin{equation}\label{eq:uv-total-majorants}
 \sum_i e_i\le|\Lambda_\ell|\delta,\qquad
 \sum_j k_j\le|\Lambda_\ell|w\,.
\end{equation}
The regular term~$\mathbf1_{\mathcal G_t(X_i)}E_i$ is bounded by~$e_i$ everywhere, so its series converges uniformly and absolutely.

\item \emph{Sum the selections.}
The disjoint-region sum is dominated by all finite selections, whose configuration-independent majorant satisfies
\begin{equation}\label{eq:uv-selection-majorant}
 \sum_{\mathcal I\subset\N\ \mathrm{finite}}\prod_{j\in\mathcal I}k_j
 =\prod_{j\ge0}(1+k_j)
 \le\exp\!\left(\sum_jk_j\right)\,.
\end{equation}
The equality follows by expanding finite products and taking increasing limits; the inequality uses~$1+s\le e^s$ for~$s\ge0$.
This proves absolute uniform convergence and the first two bounds in~\eqref{eq:uv-activity-global-bounds}.

\item \emph{Pass the properties to the sum.}
All indicators in~\eqref{eq:uv-activity-weight} are Borel and gauge invariant, as are the activities, so the limits have these properties.
The exponent in~$\rho$ is bounded using~$|S(U)|\le\|S\|_\infty$ and~$|E_t(U)|\le|\Lambda_\ell|\delta$; thus
\begin{equation*}
 |\rho(U)|\le e^{c|\Lambda_\ell|+\|S\|_\infty+|\Lambda_\ell|(\delta+w)}.
\end{equation*}
The product of the normalised Haar measures in Definition~\ref{def:wilson-law}, equation~\eqref{eq:wilson-law}, is a probability measure, so this bound gives integrability.
If there are no bad sites, a nonempty selection cannot have empty union of its nonempty~$I_j$; only the empty selection remains.
\end{enumerate}
\end{proof}

The estimates here allow the local bounds~$\delta,w$ to be independent of volume, while their global sums grow with~$|\Lambda_\ell|$.
They justify summation and integration of supplied activities.
